\documentclass[UTF8,12pt]{ctexart}
\usepackage[a4paper,margin=24mm]{geometry}
\usepackage{amsmath,amssymb,bm,graphicx,adjustbox}
\newcommand{\fitmath}[1]{\begin{adjustbox}{max width=\linewidth}$\displaystyle #1$\end{adjustbox}}
\setlength{\parindent}{0pt}
\setlength{\parskip}{.7em}
\sloppy
\allowdisplaybreaks
\begin{document}
\section*{1991 年考研数学三 · 真题}
\subsection*{1991年全国硕士研究生招生考试试题}

\subsection*{（试卷IV）}

\subsection*{一、填空题（本题共5小题，每小题3分，满分15分）}

（1）设 \(\displaystyle z=e^{\sin(xy)}\)，则 \(\displaystyle dz=\)\_\_\_\_\_\_。


（2）设曲线 \(\displaystyle f(x)=x^3+ax\) 与 \(\displaystyle g(x)=bx^2+c\) 都通过点 \(\displaystyle (-1,\allowbreak 0)\)，且在点 \(\displaystyle (-1,\allowbreak 0)\) 有公共切线，则 \(\displaystyle a=\)\_\_\_\_\_\_，\(\displaystyle b=\)\_\_\_\_\_\_，\(\displaystyle c=\)\_\_\_\_\_\_。


（3）设 \(\displaystyle f(x)=xe^x\)，则 \(\displaystyle f^{(n)}(x)\) 在点 \(\displaystyle x=\)\_\_\_\_\_\_处取极小值\_\_\_\_\_\_。


（4）设 \(\displaystyle A\) 和 \(\displaystyle B\) 为可逆矩阵，\(\displaystyle X=\begin{pmatrix}O&A\\B&O\end{pmatrix}\) 为分块矩阵，则 \(\displaystyle X^{-1}=\)\_\_\_\_\_\_。


（5）设随机变量 \(\displaystyle X\) 的分布函数为


\[\fitmath{\displaystyle F(x)=P\{X\le x\}=\begin{cases}0,&x<-1,\\0.4,&-1\le x<1,\\0.8,&1\le x<3,\\1,&x\ge3,\end{cases}}\]


则 \(\displaystyle X\) 的概率分布为\_\_\_\_\_\_。


\subsection*{二、选择题（本题共5小题，每小题3分，满分15分）}

（1）下列各式中正确的是（　　）


（A）\(\displaystyle \lim_{x\to0^+}(1+1/x)^x=1\)。 （B）\(\displaystyle \lim_{x\to0^+}(1+1/x)^x=e\)。 （C）\(\displaystyle \lim_{x\to\infty}(1-1/x)^x=-e\)。 （D）\(\displaystyle \lim_{x\to\infty}(1+1/x)^{-x}=e\)。


（2）设 \(\displaystyle 0\le a_n<\dfrac{\displaystyle 1}{\displaystyle n}\;(n=1,\allowbreak 2,\allowbreak \ldots)\)，则下列级数中肯定收敛的是（　　）


（A）\(\displaystyle \sum_{n=1}^{\infty}a_n\)。 （B）\(\displaystyle \sum_{n=1}^{\infty}(-1)^na_n\)。 （C）\(\displaystyle \sum_{n=1}^{\infty}\sqrt{a_n}\)。 （D）\(\displaystyle \sum_{n=1}^{\infty}(-1)^na_n^2\)。


（3）设 \(\displaystyle A\) 为 \(\displaystyle n\) 阶可逆矩阵，\(\displaystyle \lambda\) 是 \(\displaystyle A\) 的一个特征值，则 \(\displaystyle A\) 的伴随矩阵 \(\displaystyle A^*\) 的特征值之一是（　　）


（A）\(\displaystyle \lambda^{-1}|A|^n\)。 （B）\(\displaystyle \lambda^{-1}|A|\)。 （C）\(\displaystyle \lambda|A|\)。 （D）\(\displaystyle \lambda|A|^n\)。


（4）设 \(\displaystyle A\) 和 \(\displaystyle B\) 是任意两个概率不为零的互不相容事件，则下列结论中肯定正确的是（　　）


（A）\(\displaystyle \bar A\) 与 \(\displaystyle \bar B\) 不相容。 （B）\(\displaystyle \bar A\) 与 \(\displaystyle \bar B\) 相容。 （C）\(\displaystyle P(AB)=P(A)P(B)\)。 （D）\(\displaystyle P(A-B)=P(A)\)。


（5）对于任意两个随机变量 \(\displaystyle X\) 和 \(\displaystyle Y\)，若 \(\displaystyle E(XY)=E(X)\cdot E(Y)\)，则（　　）


（A）\(\displaystyle D(XY)=D(X)D(Y)\)。 （B）\(\displaystyle D(X+Y)=D(X)+D(Y)\)。 （C）\(\displaystyle X\) 与 \(\displaystyle Y\) 独立。 （D）\(\displaystyle X\) 与 \(\displaystyle Y\) 不独立。


\subsection*{三、（本题满分5分）}

求极限 \(\displaystyle \lim_{x\to0}\left(\dfrac{\displaystyle e^x+e^{2x}+\cdots+e^{nx}}{\displaystyle n}\right)^{1/x}\)，其中 \(\displaystyle n\) 是给定的自然数。


\subsection*{四、（本题满分5分）}

计算二重积分 \(\displaystyle I=\displaystyle \iint_D y\,dx\,dy\)，其中 \(\displaystyle D\) 是由 \(\displaystyle x\) 轴、\(\displaystyle y\) 轴与曲线 \(\displaystyle \sqrt{x/a}+\sqrt{y/b}=1\) 所围成的区域，\(\displaystyle a>0,\allowbreak b>0\)。


\subsection*{五、（本题满分5分）}

求微分方程 \(\displaystyle xy\dfrac{\displaystyle dy}{\displaystyle dx}=x^2+y^2\) 满足条件 \(\displaystyle y|_{x=e}=2e\) 的特解。


\subsection*{六、（本题满分5分）}

假设曲线 \(\displaystyle L_1:y=1-x^2\;(0\le x\le1)\)、\(\displaystyle x\) 轴和 \(\displaystyle y\) 轴所围区域被曲线 \(\displaystyle L_2:y=ax^2\) 分为面积相等的两部分，其中 \(\displaystyle a\) 是大于零的常数，试确定 \(\displaystyle a\) 的值。


\subsection*{七、（本题满分8分）}

某厂家生产的一种产品同时在两个市场销售，售价分别为 \(\displaystyle p_1\) 和 \(\displaystyle p_2\)，销售量分别为 \(\displaystyle q_1\) 和 \(\displaystyle q_2\)，需求函数分别为 \(\displaystyle q_1=24-0.2p_1\) 和 \(\displaystyle q_2=10-0.05p_2\)，总成本函数为 \(\displaystyle C=35+40(q_1+q_2)\)。试问：厂家如何确定两个市场的售价，能使其获得的总利润最大？最大总利润为多少？


\subsection*{八、（本题满分6分）}

试证明函数 \(\displaystyle f(x)=(1+\dfrac{\displaystyle 1}{\displaystyle x})^x\) 在区间 \(\displaystyle (0,\allowbreak +\infty)\) 内单调增加。


\subsection*{九、（本题满分7分）}

设有3维列向量


\[\fitmath{\displaystyle \boldsymbol\alpha_1=\begin{pmatrix}1+\lambda\\1\\1\end{pmatrix},\quad\boldsymbol\alpha_2=\begin{pmatrix}1\\1+\lambda\\1\end{pmatrix},\quad\boldsymbol\alpha_3=\begin{pmatrix}1\\1\\1+\lambda\end{pmatrix},\quad\boldsymbol\beta=\begin{pmatrix}0\\\lambda\\\lambda^2\end{pmatrix}.}\]


问 \(\displaystyle \lambda\) 取何值时，（1）\(\displaystyle \boldsymbol\beta\) 可由 \(\displaystyle \boldsymbol\alpha_1,\allowbreak \boldsymbol\alpha_2,\allowbreak \boldsymbol\alpha_3\) 线性表示，且表达式唯一？（2）可以线性表示，但表达式不唯一？（3）不能线性表示？


\subsection*{十、（本题满分6分）}

考虑二次型 \(\displaystyle f=x_1^2+4x_2^2+4x_3^2+2\lambda x_1x_2-2x_1x_3+4x_2x_3\)，问 \(\displaystyle \lambda\) 取何值时，\(\displaystyle f\) 为正定二次型？


\subsection*{十一、（本题满分6分）}

试证明 \(\displaystyle n\) 维列向量组 \(\displaystyle \boldsymbol\alpha_1,\allowbreak \boldsymbol\alpha_2,\allowbreak \ldots,\allowbreak \boldsymbol\alpha_n\) 线性无关的充分必要条件是


\[\fitmath{\displaystyle D=\begin{vmatrix}\boldsymbol\alpha_1^{\mathrm T}\boldsymbol\alpha_1&\boldsymbol\alpha_1^{\mathrm T}\boldsymbol\alpha_2&\cdots&\boldsymbol\alpha_1^{\mathrm T}\boldsymbol\alpha_n\\\boldsymbol\alpha_2^{\mathrm T}\boldsymbol\alpha_1&\boldsymbol\alpha_2^{\mathrm T}\boldsymbol\alpha_2&\cdots&\boldsymbol\alpha_2^{\mathrm T}\boldsymbol\alpha_n\\\vdots&\vdots&\ddots&\vdots\\\boldsymbol\alpha_n^{\mathrm T}\boldsymbol\alpha_1&\boldsymbol\alpha_n^{\mathrm T}\boldsymbol\alpha_2&\cdots&\boldsymbol\alpha_n^{\mathrm T}\boldsymbol\alpha_n\end{vmatrix}\ne0,}\]


其中 \(\displaystyle \boldsymbol\alpha_i^{\mathrm T}\) 表示列向量 \(\displaystyle \boldsymbol\alpha_i\) 的转置，\(\displaystyle i=1,\allowbreak 2,\allowbreak \ldots,\allowbreak n\)。


\subsection*{十二、（本题满分6分）}

一汽车沿一街道行驶，需要通过三个均设有红绿信号灯的路口，每个信号灯为红或绿与其他信号灯为红或绿相互独立，且红绿两种信号显示的时间相等。以 \(\displaystyle X\) 表示该汽车首次遇到红灯前已通过的路口的个数。求 \(\displaystyle X\) 的概率分布。


\subsection*{十三、（本题满分6分）}

假设随机变量 \(\displaystyle X\) 和 \(\displaystyle Y\) 在圆域 \(\displaystyle x^2+y^2\le r^2\) 上服从联合均匀分布。（1）求 \(\displaystyle X\) 和 \(\displaystyle Y\) 的相关系数 \(\displaystyle \rho\)；（2）问 \(\displaystyle X\) 和 \(\displaystyle Y\) 是否独立？


\subsection*{十四、（本题满分5分）}

设总体 \(\displaystyle X\) 的概率密度为


\[\fitmath{\displaystyle p(x,\lambda)=\begin{cases}\lambda a x^{a-1}e^{-\lambda x^a},&x>0,\\0,&x\le0,\end{cases}}\]


其中 \(\displaystyle \lambda>0\) 是未知参数，\(\displaystyle a>0\) 是已知常数。试根据来自总体 \(\displaystyle X\) 的简单随机样本 \(\displaystyle X_1,\allowbreak X_2,\allowbreak \ldots,\allowbreak X_n\)，求 \(\displaystyle \lambda\) 的最大似然估计量 \(\displaystyle \hat\lambda\)。


\subsection*{（试卷V）}

\subsection*{一、填空题（本题共5小题，每小题3分，满分15分）}

（1）【同试卷IV 第一、（1）题】 （2）【同试卷IV 第一、（2）题】 （3）【同试卷IV 第一、（3）题】


（4）\(\displaystyle n\) 阶行列式


\[\fitmath{\displaystyle \begin{vmatrix}a&b&0&\cdots&0&0\\0&a&b&\cdots&0&0\\0&0&a&\cdots&0&0\\\vdots&\vdots&\vdots&\ddots&\vdots&\vdots\\0&0&0&\cdots&a&b\\b&0&0&\cdots&0&a\end{vmatrix}=\underline{\qquad}.}\]


（5）设 \(\displaystyle A,\allowbreak B\) 为随机事件，\(\displaystyle P(A)=0.7,\allowbreak P(A-B)=0.3\)，则 \(\displaystyle P(\overline{AB})=\)\_\_\_\_\_\_。


\subsection*{二、选择题（本题共5小题，每小题3分，满分15分）}

（1）【同试卷IV 第二、（1）题】


（2）设数列的通项为 \(\displaystyle x_n=\begin{cases}\dfrac{\displaystyle n^2+\sqrt n}{\displaystyle n},\allowbreak &n\text{ 为奇数},\allowbreak \\\dfrac{\displaystyle 1}{\displaystyle n},\allowbreak &n\text{ 为偶数},\allowbreak \end{cases}\) 则当 \(\displaystyle n\to\infty\) 时，\(\displaystyle x_n\) 是（　　） （A）无穷大量。 （B）无穷小量。 （C）有界变量。 （D）无界变量。


（3）设 \(\displaystyle A\) 与 \(\displaystyle B\) 为 \(\displaystyle n\) 阶方阵，且 \(\displaystyle AB=O\)，则必有（　　） （A）\(\displaystyle A=O\) 或 \(\displaystyle B=O\)。 （B）\(\displaystyle AB=BA\)。 （C）\(\displaystyle |A|=0\) 或 \(\displaystyle |B|=0\)。 （D）\(\displaystyle |A|+|B|=0\)。


（4）设 \(\displaystyle A\) 是 \(\displaystyle m\times n\) 矩阵，\(\displaystyle A\boldsymbol x=0\) 是非齐次线性方程组 \(\displaystyle A\boldsymbol x=\boldsymbol b\) 所对应的齐次线性方程组，则下列结论正确的是（　　）


（A）若 \(\displaystyle A\boldsymbol x=0\) 仅有零解，则 \(\displaystyle A\boldsymbol x=\boldsymbol b\) 有唯一解。 （B）若 \(\displaystyle A\boldsymbol x=0\) 有非零解，则 \(\displaystyle A\boldsymbol x=\boldsymbol b\) 有无穷多个解。 （C）若 \(\displaystyle A\boldsymbol x=\boldsymbol b\) 有无穷多个解，则 \(\displaystyle A\boldsymbol x=0\) 仅有零解。 （D）若 \(\displaystyle A\boldsymbol x=\boldsymbol b\) 有无穷多个解，则 \(\displaystyle A\boldsymbol x=0\) 有非零解。


（5）【同试卷IV 第二、（4）题】


\subsection*{三、（本题满分5分）}

求极限 \(\displaystyle \lim_{x\to+\infty}(x+\sqrt{1+x^2})^{1/x}\)。


\subsection*{四、（本题满分5分）}

求定积分 \(\displaystyle I=\displaystyle \int_{-1}^1(2x+|x|+1)^2\,dx\)。


\subsection*{五、（本题满分5分）}

求不定积分 \(\displaystyle I=\displaystyle \int\dfrac{\displaystyle x^2}{\displaystyle 1+x^2}\arctan x\,dx\)。


\subsection*{六、（本题满分5分）}

已知 \(\displaystyle xy=xf(z)+yg(z)\)，\(\displaystyle xf\prime(z)+yg\prime(z)\ne0\)，其中 \(\displaystyle z=z(x,\allowbreak y)\) 是 \(\displaystyle x\) 和 \(\displaystyle y\) 的函数，求证 \(\displaystyle [x-g(z)]\dfrac{\displaystyle \partial z}{\displaystyle \partial x}=[y-f(z)]\dfrac{\displaystyle \partial z}{\displaystyle \partial y}\)。


\subsection*{七、（本题满分6分）【同试卷IV 第六题】}

\subsection*{八、（本题满分8分）【同试卷IV 第七题】}

\subsection*{九、（本题满分6分）}

证明不等式 \(\displaystyle \ln(1+\dfrac{\displaystyle 1}{\displaystyle x})>\dfrac{\displaystyle 1}{\displaystyle 1+x}\;(0<x<+\infty)\)。


\subsection*{十、（本题满分5分）}

设 \(\displaystyle n\) 阶矩阵 \(\displaystyle A\) 和 \(\displaystyle B\) 满足条件 \(\displaystyle A+B=AB\)。（1）证明 \(\displaystyle A-E\) 为可逆矩阵；（2）已知


\[\fitmath{\displaystyle B=\begin{pmatrix}1&-3&0\\2&1&0\\0&0&2\end{pmatrix},}\]


求矩阵 \(\displaystyle A\)。


\subsection*{十一、（本题满分7分）【同试卷IV 第九题】}

\subsection*{十二、（本题满分5分）}

已知向量 \(\displaystyle \boldsymbol\alpha=(1,\allowbreak k,\allowbreak 1)^{\mathrm T}\) 是矩阵 \(\displaystyle A=\begin{pmatrix}2&1&1\\1&2&1\\1&1&2\end{pmatrix}\) 的逆矩阵 \(\displaystyle A^{-1}\) 的特征向量，试求常数 \(\displaystyle k\) 的值。


\subsection*{十三、（本题满分7分）}

一汽车沿一街道行驶，需要通过三个均设有红绿信号灯的路口，每个信号灯为红或绿与其他信号灯为红或绿相互独立，且红绿两种信号显示的时间相等。以 \(\displaystyle X\) 表示该汽车首次遇到红灯前已通过的路口的个数。（1）求 \(\displaystyle X\) 的概率分布；（2）求 \(\displaystyle E(\dfrac{\displaystyle 1}{\displaystyle 1+X})\)。


\subsection*{十四、（本题满分6分）}

在电源电压不超过200伏、200～240伏和超过240伏三种情况下，某种电子元件损坏的概率分别为0.1、0.001和0.2。假设电源电压 \(\displaystyle X\) 服从正态分布 \(\displaystyle N(220,\allowbreak 25^2)\)。试求：（1）该电子元件损坏的概率 \(\displaystyle \alpha\)；（2）该电子元件损坏时，电源电压在220～240伏的概率 \(\displaystyle \beta\)。附表（表中 \(\displaystyle \Phi(x)\) 是标准正态分布函数）


\[\fitmath{\displaystyle \begin{array}{c|cccccccc}x&0.10&0.20&0.40&0.60&0.80&1.00&1.20&1.40\\\hline\Phi(x)&0.530&0.579&0.655&0.726&0.788&0.841&0.885&0.919\end{array}}\]

\end{document}
